On écarte le solide $\mathrm{(S)}$ (Fig.~\ref{fig:solide_ressort}) de sa
position d'équilibre et on le libère sans vitesse initiale à l'instant
$t_{0}=0$.

\begin{minipage}{0.70\linewidth}
    \begin{donnees}
        \begin{itemize}
          \item Tous les frottements sont négligeables~;
          \item On choisit l'état où le ressort n'est pas déformé comme
                référence de l'énergie potentielle élastique $E_{pe}$ et le plan
                horizontal contenant $G$ comme état de référence de l'énergie
                potentielle de pesanteur $E_{pp}$.
        \end{itemize}
    \end{donnees}

    La courbe de la Fig.~\ref{fig:2bac_svt_national_2016_normale_phys_ex3_3}
    représente les variations de $E_{pe}$ en fonction de $x^{2}$, carré de
    l'abscisse $x$ du centre d'inertie $G$ dans le repère $(O,\ex)$.
    \begin{enumerate}
        \item En exploitant la courbe de la
              Fig.~\ref{fig:2bac_svt_national_2016_normale_phys_ex3_3},
              trouver les valeurs de~:
              \begin{enumerate}
                  \item la constante de raideur $K$.
                  \item l'énergie potentielle élastique maximale $E_{pe,\max}$.
                  \item l'amplitude $X_{m}$ des oscillations.
              \end{enumerate}
        \item Déduire, en justifiant votre réponse, la valeur de l'énergie
              mécanique $E_{m}$ du système oscillant.
    \end{enumerate}
\end{minipage}
\hfill
\begin{minipage}{0.29\linewidth}
    \begin{figure}[H]
        \centering
        \begin{tikzpicture}
            \def\K{1}
            \begin{axis}[
                    xmin=0,xmax=19,
                    ymin=0,ymax=9.8,
                    width=6.5cm,height=6.5cm,
                    xlabel=$x^{2}~(\times\qty{1e-4}{\square\meter})$,
                    ylabel=$E_{pe}~(\unit{\mJ})$,
                    scaled ticks=false,
                    xtick distance=4, ytick distance=2,
                    xticklabels={,,4}, yticklabels={,,2},
                    minor tick num=0,
                ]
                \addplot[domain=0:16,samples=200,
                    very thick,blue] {0.5*\K*x}
                    coordinate[pos=1](A);
            \end{axis}
            \draw[very thick] ([yshift=-5mm]A) -- ([yshift=5mm]A);
            \foreach \f in {0.02,0.14,...,1}
                \draw[thick]
                    ($([yshift=-5mm]A)!\f!([yshift=5mm]A)$) -- ++(-20:.2);
        \end{tikzpicture}
        \caption{}
        \label{fig:2bac_svt_national_2016_normale_phys_ex3_3}
    \end{figure}
\end{minipage}

\begin{enumerate}
    \setcounter{enumi}{2}
    \item Le centre d'inertie $G$ passe par la position d'équilibre dans le sens
          positif avec la vitesse $v=\qty{0,25}{\v}$.

          Montrer que l'expression de la période propre des oscillations s'écrit
          $T_{0}=2\pi\cdot\frac{X_{m}}{v}$. Calculer $T_{0}$.
\end{enumerate}

